%% tuftedarklayout.tex -- the improved tuftedark layout, as worked out in the
%% Intelligence Engineering deck, for every other deck to share:
%%   \usetheme{tuftedark}
%%   %% tuftedarklayout.tex -- the improved tuftedark layout, as worked out in the
%% Intelligence Engineering deck, for every other deck to share:
%%   \usetheme{tuftedark}
%%   %% tuftedarklayout.tex -- the improved tuftedark layout, as worked out in the
%% Intelligence Engineering deck, for every other deck to share:
%%   \usetheme{tuftedark}
%%   %% tuftedarklayout.tex -- the improved tuftedark layout, as worked out in the
%% Intelligence Engineering deck, for every other deck to share:
%%   \usetheme{tuftedark}
%%   \input{tuftedarklayout}
%% (sharedAssets/beamer is on TEXINPUTS through each deck's .latexmkrc.)
%%
%% What it changes against the bare theme:
%%   footer       speaker name, section, progress track, counter; no dots
%%   \sidecaption up the strip's inner edge in typewriter grey, off the image
%%   statement, takeaway   red rule, flush left, bold mixed case, \alertbf
%%   \framesubtitle        printed, grey, above the title rule
%%   \titleontop           redraws a long title over a side strip
%%
%% COPIED VERBATIM from notes_intelligenceengineering/slides/
%% intelligence_engineering.tex (2026-11-20), which still carries its own
%% copy. The comments speak of "this deck" for that reason. Change both, or
%% move that deck onto this file.

\makeatletter
% \sidecaption shadowed for this deck. The theme sets the caption horizontal,
% white and italic INSIDE the strip, which a photo swallows; here it reads up
% the strip's outer edge in the placeholder's own typewriter grey, mirroring
% \slidesources on the opposite page edge, so a real image and its
% \phimgside placeholder caption the same way. 13mm clears the footer band.
% The node is filled with the page's own card colour and given a little
% horizontal padding, so wherever a long caption runs up past the frametitle's
% hairline rule it masks that rule instead of crossing it. The caption always
% sits OUTSIDE the image strip, on the card, so the fill is invisible; it only
% ever paints over the rule. Keep captions under ~50 characters anyway: the
% node has no text width and the page clips whatever runs past the top edge.
\renewcommand{\sidecaption}[1]{%
  \setlength{\td@stripw}{\td@sidefrac\paperwidth}%
  \def\td@westtest{west}%
  \begin{tikzpicture}[remember picture,overlay]
    \ifx\td@sideside\td@westtest
      \node[rotate=90,anchor=west,inner xsep=1.2mm,inner ysep=0.6mm,
            fill=tdcard,text=tddim,font=\ttfamily\scriptsize]
        at ([xshift=\dimexpr\td@stripw+2mm\relax,yshift=13mm]current page.south west){#1};
    \else
      \node[rotate=90,anchor=west,inner xsep=1.2mm,inner ysep=0.6mm,
            fill=tdcard,text=tddim,font=\ttfamily\scriptsize]
        at ([xshift=\dimexpr-\td@stripw-2mm\relax,yshift=13mm]current page.south east){#1};
    \fi
  \end{tikzpicture}%
}
\makeatother

% A References frame that runs long is set [allowframebreaks]; beamer then
% appends a roman continuation counter to the title ("References II"), which
% reads like a second section. Blank it: both pages are just References.
\setbeamertemplate{frametitle continuation}{}

% ------------------------------------------------------------- statement
% The punchline slide, overridden deck level. The theme's own statement is
% letterspaced uppercase with a red uppercase second line; both lines shout and
% the red one competes with the first instead of supporting it.
%
% This is the Disciplined Vibe closing card's idea, editorial rather than
% centred: a short red rule, the claim flush left in large bold MIXED case with
% the red carried inline by \alertbf, then the consequence one step down in
% grey. The rule is the same motif as the subchapter divider, so dividers and
% punchlines read as one family, and flush left means a wrapping claim breaks
% cleanly instead of stranding an orphan word in the middle of the slide.
%
%   \begin{statement}
%     The algorithm is the \alertbf{small box in the middle}
%     \substatement{Everything else is this course.}
%   \end{statement}
%
% The theme's versions are deliberately shadowed, not patched: other decks share
% that file. \RenewDocumentEnvironment/Command so this fails loudly if the theme
% ever stops providing them.
\RenewDocumentEnvironment{statement}{+b}{%
  \par\vfill
  {\color{DHBWred}\rule{0.10\paperwidth}{2.5pt}}\par\vskip1.1em
  {\fontsize{28}{33}\selectfont\bfseries\raggedright #1\par}%
  \vfill\par}{}
% \substatement stays \protected so the surrounding grab cannot expand it.
\RenewDocumentCommand{\substatement}{m}{%
  \par\vskip0.9em{\large\color{tddim}\raggedright #1\par}}

% -------------------------------------------------------------- takeaway
% Also overridden deck level, and for the same reason as statement above: the
% theme's takeaway is \Large italic red centred, which shouts in a typeface
% that was never drawn to shout, and centring strands an orphan word when the
% line wraps. Here it joins the statement family instead: red rule, claim flush
% left in bold MIXED case, the red carried inline by \alertbf rather than
% painted over the whole sentence. One step smaller than statement, because a
% takeaway closes a chapter where a statement opens a question.
%
%   \begin{takeaway}
%     The gap between the two is the \alertbf{machine learning system}.
%   \end{takeaway}
\RenewDocumentEnvironment{takeaway}{+b}{%
  \par\vfill
  {\color{DHBWred}\rule{0.07\paperwidth}{2.5pt}}\par\vskip0.9em
  {\fontsize{22}{27}\selectfont\bfseries\raggedright #1\par}%
  \vfill\par}{}

% ---------------------------------------------------------- frame subtitle
% The theme's frametitle template prints the title and the hairline rule and
% drops \insertframesubtitle on the floor, so \framesubtitle{...} renders
% nothing. It is reinstated here, deck level, with the theme's own template
% copied verbatim and one branch added: grey, small, above the rule.
% Used for the source's own one line under a heading, e.g. the condition that
% names each of the five keyphrases in chapter 1.
\setbeamerfont{framesubtitle}{size=\small}
\setbeamercolor{framesubtitle}{fg=tddim}
\makeatletter
\setbeamertemplate{frametitle}{%
  \nointerlineskip%
  \vskip2.4ex%
  {\usebeamerfont{frametitle}\usebeamercolor[fg]{frametitle}\strut\tikz[remember picture,overlay]\coordinate (tdframetitle);\insertframetitle\par}%
  \ifx\insertframesubtitle\@empty\else
    \vskip0.35ex%
    {\usebeamerfont{framesubtitle}\usebeamercolor[fg]{framesubtitle}%
      \insertframesubtitle\par}%
  \fi
  \vskip0.2ex%
  {\color{tdrule}\rule{\textwidth}{0.4pt}}%
  \vskip0.4ex%
}
\makeatother

% Title on top of a side strip: \titleontop, called after \sideimagecover.
% The frametitle is typeset before the frame body, so a strip laid down in
% the body covers a long title's right end. The template above marks where
% the title starts with a zero-size coordinate; this draws the title again
% at that mark, over the photo.
\newcommand{\titleontop}{%
  \begin{tikzpicture}[remember picture,overlay]
    \node[anchor=base west,inner sep=0pt] at (tdframetitle)
      {\usebeamerfont{frametitle}\usebeamercolor[fg]{frametitle}\insertframetitle};
  \end{tikzpicture}%
}

% ------------------------------------------------------------- footer nav
% The theme's footer is beamer's mini-frame navigation: one hollow dot per
% frame, grouped under its section name. It is sized for a deck the shape of
% Disciplined Vibe, eight sections over fifty frames. This deck mirrors a
% seventeen block course: sixteen sections over about a hundred and eighty
% frames, which overruns the footer band by some 340pt on every single page.
% Neither knob fixes it, the dots alone are already three times the paper
% width. So the nav is replaced here, at deck level, with the same information
% compressed: the section you are in on the left, a progress track in the
% middle, the counter on the right. The theme is untouched and the other decks
% keep their dots. The track is drawn in tdfg on a tdrule ground rather than in
% red, so the footer never spends the slide's one red element.
% Drawn with \rule rather than tikz on purpose: a tikz coordinate resolves a
% bare length register through pgfmath, which hands back a unitless number that
% tikz then reads as centimetres, and the bar blows past the maximum dimension.
\newlength{\ieprogwd}
\newlength{\ieprogfill}
\setlength{\ieprogwd}{0.30\paperwidth}
% The fraction goes through pgfmath, which works in floating point. Doing it in
% TeX dimen arithmetic does not survive the first pass: before the aux file is
% written \inserttotalframenumber is stale, the \divide leaves \ieprogwd whole,
% and \ieprogwd times a frame number past about 120 exceeds the maximum TeX
% dimension. The clamp covers the same stale-aux case, where the current frame
% number can exceed the recorded total and the fill would run past its track.
\newcommand{\ieprogressbar}{%
  \ifnum\inserttotalframenumber>0\relax
    \ifnum\insertframenumber<\inserttotalframenumber\relax
      % The ratio is parenthesised so it is evaluated first. pgfmath keeps its
      % intermediates in a dimen register, so multiplying the track width by the
      % frame number before dividing overflows once past frame 120.
      \pgfmathsetlength{\ieprogfill}{\ieprogwd*(\insertframenumber/\inserttotalframenumber)}%
    \else
      \setlength{\ieprogfill}{\ieprogwd}%
    \fi
  \else
    \setlength{\ieprogfill}{0pt}%
  \fi
  \raisebox{0.6ex}{\makebox[\ieprogwd][l]{%
    \rlap{\textcolor{tdrule}{\rule{\ieprogwd}{1.1pt}}}%
    \textcolor{tdfg}{\rule{\ieprogfill}{1.1pt}}}}%
}
% The speaker's name leads the footer on every slide, so a photo taken during
% the talk carries it (Mark, 2026-10-26). The right edge watermark cannot do
% this: side strips and full page sketches cover it.
\newcommand{\iefootname}{Prof.\ Dr.-Ing.\ Mark Schutera\enspace\textperiodcentered\enspace}
\setbeamertemplate{footline}{%
  \begin{beamercolorbox}[wd=\paperwidth,ht=6mm,dp=5mm,leftskip=6mm,rightskip=6mm]{section in head/foot}%
    {\usebeamerfont{section in head/foot}\color{tddim}\iefootname\insertsectionhead}%
    \hfill
    \ieprogressbar
    \hfill
    {\usebeamerfont{section in head/foot}\color{tddim}%
      \insertframenumber\,/\,\inserttotalframenumber}%
  \end{beamercolorbox}%
}

%% (sharedAssets/beamer is on TEXINPUTS through each deck's .latexmkrc.)
%%
%% What it changes against the bare theme:
%%   footer       speaker name, section, progress track, counter; no dots
%%   \sidecaption up the strip's inner edge in typewriter grey, off the image
%%   statement, takeaway   red rule, flush left, bold mixed case, \alertbf
%%   \framesubtitle        printed, grey, above the title rule
%%   \titleontop           redraws a long title over a side strip
%%
%% COPIED VERBATIM from notes_intelligenceengineering/slides/
%% intelligence_engineering.tex (2026-11-20), which still carries its own
%% copy. The comments speak of "this deck" for that reason. Change both, or
%% move that deck onto this file.

\makeatletter
% \sidecaption shadowed for this deck. The theme sets the caption horizontal,
% white and italic INSIDE the strip, which a photo swallows; here it reads up
% the strip's outer edge in the placeholder's own typewriter grey, mirroring
% \slidesources on the opposite page edge, so a real image and its
% \phimgside placeholder caption the same way. 13mm clears the footer band.
% The node is filled with the page's own card colour and given a little
% horizontal padding, so wherever a long caption runs up past the frametitle's
% hairline rule it masks that rule instead of crossing it. The caption always
% sits OUTSIDE the image strip, on the card, so the fill is invisible; it only
% ever paints over the rule. Keep captions under ~50 characters anyway: the
% node has no text width and the page clips whatever runs past the top edge.
\renewcommand{\sidecaption}[1]{%
  \setlength{\td@stripw}{\td@sidefrac\paperwidth}%
  \def\td@westtest{west}%
  \begin{tikzpicture}[remember picture,overlay]
    \ifx\td@sideside\td@westtest
      \node[rotate=90,anchor=west,inner xsep=1.2mm,inner ysep=0.6mm,
            fill=tdcard,text=tddim,font=\ttfamily\scriptsize]
        at ([xshift=\dimexpr\td@stripw+2mm\relax,yshift=13mm]current page.south west){#1};
    \else
      \node[rotate=90,anchor=west,inner xsep=1.2mm,inner ysep=0.6mm,
            fill=tdcard,text=tddim,font=\ttfamily\scriptsize]
        at ([xshift=\dimexpr-\td@stripw-2mm\relax,yshift=13mm]current page.south east){#1};
    \fi
  \end{tikzpicture}%
}
\makeatother

% A References frame that runs long is set [allowframebreaks]; beamer then
% appends a roman continuation counter to the title ("References II"), which
% reads like a second section. Blank it: both pages are just References.
\setbeamertemplate{frametitle continuation}{}

% ------------------------------------------------------------- statement
% The punchline slide, overridden deck level. The theme's own statement is
% letterspaced uppercase with a red uppercase second line; both lines shout and
% the red one competes with the first instead of supporting it.
%
% This is the Disciplined Vibe closing card's idea, editorial rather than
% centred: a short red rule, the claim flush left in large bold MIXED case with
% the red carried inline by \alertbf, then the consequence one step down in
% grey. The rule is the same motif as the subchapter divider, so dividers and
% punchlines read as one family, and flush left means a wrapping claim breaks
% cleanly instead of stranding an orphan word in the middle of the slide.
%
%   \begin{statement}
%     The algorithm is the \alertbf{small box in the middle}
%     \substatement{Everything else is this course.}
%   \end{statement}
%
% The theme's versions are deliberately shadowed, not patched: other decks share
% that file. \RenewDocumentEnvironment/Command so this fails loudly if the theme
% ever stops providing them.
\RenewDocumentEnvironment{statement}{+b}{%
  \par\vfill
  {\color{DHBWred}\rule{0.10\paperwidth}{2.5pt}}\par\vskip1.1em
  {\fontsize{28}{33}\selectfont\bfseries\raggedright #1\par}%
  \vfill\par}{}
% \substatement stays \protected so the surrounding grab cannot expand it.
\RenewDocumentCommand{\substatement}{m}{%
  \par\vskip0.9em{\large\color{tddim}\raggedright #1\par}}

% -------------------------------------------------------------- takeaway
% Also overridden deck level, and for the same reason as statement above: the
% theme's takeaway is \Large italic red centred, which shouts in a typeface
% that was never drawn to shout, and centring strands an orphan word when the
% line wraps. Here it joins the statement family instead: red rule, claim flush
% left in bold MIXED case, the red carried inline by \alertbf rather than
% painted over the whole sentence. One step smaller than statement, because a
% takeaway closes a chapter where a statement opens a question.
%
%   \begin{takeaway}
%     The gap between the two is the \alertbf{machine learning system}.
%   \end{takeaway}
\RenewDocumentEnvironment{takeaway}{+b}{%
  \par\vfill
  {\color{DHBWred}\rule{0.07\paperwidth}{2.5pt}}\par\vskip0.9em
  {\fontsize{22}{27}\selectfont\bfseries\raggedright #1\par}%
  \vfill\par}{}

% ---------------------------------------------------------- frame subtitle
% The theme's frametitle template prints the title and the hairline rule and
% drops \insertframesubtitle on the floor, so \framesubtitle{...} renders
% nothing. It is reinstated here, deck level, with the theme's own template
% copied verbatim and one branch added: grey, small, above the rule.
% Used for the source's own one line under a heading, e.g. the condition that
% names each of the five keyphrases in chapter 1.
\setbeamerfont{framesubtitle}{size=\small}
\setbeamercolor{framesubtitle}{fg=tddim}
\makeatletter
\setbeamertemplate{frametitle}{%
  \nointerlineskip%
  \vskip2.4ex%
  {\usebeamerfont{frametitle}\usebeamercolor[fg]{frametitle}\strut\tikz[remember picture,overlay]\coordinate (tdframetitle);\insertframetitle\par}%
  \ifx\insertframesubtitle\@empty\else
    \vskip0.35ex%
    {\usebeamerfont{framesubtitle}\usebeamercolor[fg]{framesubtitle}%
      \insertframesubtitle\par}%
  \fi
  \vskip0.2ex%
  {\color{tdrule}\rule{\textwidth}{0.4pt}}%
  \vskip0.4ex%
}
\makeatother

% Title on top of a side strip: \titleontop, called after \sideimagecover.
% The frametitle is typeset before the frame body, so a strip laid down in
% the body covers a long title's right end. The template above marks where
% the title starts with a zero-size coordinate; this draws the title again
% at that mark, over the photo.
\newcommand{\titleontop}{%
  \begin{tikzpicture}[remember picture,overlay]
    \node[anchor=base west,inner sep=0pt] at (tdframetitle)
      {\usebeamerfont{frametitle}\usebeamercolor[fg]{frametitle}\insertframetitle};
  \end{tikzpicture}%
}

% ------------------------------------------------------------- footer nav
% The theme's footer is beamer's mini-frame navigation: one hollow dot per
% frame, grouped under its section name. It is sized for a deck the shape of
% Disciplined Vibe, eight sections over fifty frames. This deck mirrors a
% seventeen block course: sixteen sections over about a hundred and eighty
% frames, which overruns the footer band by some 340pt on every single page.
% Neither knob fixes it, the dots alone are already three times the paper
% width. So the nav is replaced here, at deck level, with the same information
% compressed: the section you are in on the left, a progress track in the
% middle, the counter on the right. The theme is untouched and the other decks
% keep their dots. The track is drawn in tdfg on a tdrule ground rather than in
% red, so the footer never spends the slide's one red element.
% Drawn with \rule rather than tikz on purpose: a tikz coordinate resolves a
% bare length register through pgfmath, which hands back a unitless number that
% tikz then reads as centimetres, and the bar blows past the maximum dimension.
\newlength{\ieprogwd}
\newlength{\ieprogfill}
\setlength{\ieprogwd}{0.30\paperwidth}
% The fraction goes through pgfmath, which works in floating point. Doing it in
% TeX dimen arithmetic does not survive the first pass: before the aux file is
% written \inserttotalframenumber is stale, the \divide leaves \ieprogwd whole,
% and \ieprogwd times a frame number past about 120 exceeds the maximum TeX
% dimension. The clamp covers the same stale-aux case, where the current frame
% number can exceed the recorded total and the fill would run past its track.
\newcommand{\ieprogressbar}{%
  \ifnum\inserttotalframenumber>0\relax
    \ifnum\insertframenumber<\inserttotalframenumber\relax
      % The ratio is parenthesised so it is evaluated first. pgfmath keeps its
      % intermediates in a dimen register, so multiplying the track width by the
      % frame number before dividing overflows once past frame 120.
      \pgfmathsetlength{\ieprogfill}{\ieprogwd*(\insertframenumber/\inserttotalframenumber)}%
    \else
      \setlength{\ieprogfill}{\ieprogwd}%
    \fi
  \else
    \setlength{\ieprogfill}{0pt}%
  \fi
  \raisebox{0.6ex}{\makebox[\ieprogwd][l]{%
    \rlap{\textcolor{tdrule}{\rule{\ieprogwd}{1.1pt}}}%
    \textcolor{tdfg}{\rule{\ieprogfill}{1.1pt}}}}%
}
% The speaker's name leads the footer on every slide, so a photo taken during
% the talk carries it (Mark, 2026-10-26). The right edge watermark cannot do
% this: side strips and full page sketches cover it.
\newcommand{\iefootname}{Prof.\ Dr.-Ing.\ Mark Schutera\enspace\textperiodcentered\enspace}
\setbeamertemplate{footline}{%
  \begin{beamercolorbox}[wd=\paperwidth,ht=6mm,dp=5mm,leftskip=6mm,rightskip=6mm]{section in head/foot}%
    {\usebeamerfont{section in head/foot}\color{tddim}\iefootname\insertsectionhead}%
    \hfill
    \ieprogressbar
    \hfill
    {\usebeamerfont{section in head/foot}\color{tddim}%
      \insertframenumber\,/\,\inserttotalframenumber}%
  \end{beamercolorbox}%
}

%% (sharedAssets/beamer is on TEXINPUTS through each deck's .latexmkrc.)
%%
%% What it changes against the bare theme:
%%   footer       speaker name, section, progress track, counter; no dots
%%   \sidecaption up the strip's inner edge in typewriter grey, off the image
%%   statement, takeaway   red rule, flush left, bold mixed case, \alertbf
%%   \framesubtitle        printed, grey, above the title rule
%%   \titleontop           redraws a long title over a side strip
%%
%% COPIED VERBATIM from notes_intelligenceengineering/slides/
%% intelligence_engineering.tex (2026-11-20), which still carries its own
%% copy. The comments speak of "this deck" for that reason. Change both, or
%% move that deck onto this file.

\makeatletter
% \sidecaption shadowed for this deck. The theme sets the caption horizontal,
% white and italic INSIDE the strip, which a photo swallows; here it reads up
% the strip's outer edge in the placeholder's own typewriter grey, mirroring
% \slidesources on the opposite page edge, so a real image and its
% \phimgside placeholder caption the same way. 13mm clears the footer band.
% The node is filled with the page's own card colour and given a little
% horizontal padding, so wherever a long caption runs up past the frametitle's
% hairline rule it masks that rule instead of crossing it. The caption always
% sits OUTSIDE the image strip, on the card, so the fill is invisible; it only
% ever paints over the rule. Keep captions under ~50 characters anyway: the
% node has no text width and the page clips whatever runs past the top edge.
\renewcommand{\sidecaption}[1]{%
  \setlength{\td@stripw}{\td@sidefrac\paperwidth}%
  \def\td@westtest{west}%
  \begin{tikzpicture}[remember picture,overlay]
    \ifx\td@sideside\td@westtest
      \node[rotate=90,anchor=west,inner xsep=1.2mm,inner ysep=0.6mm,
            fill=tdcard,text=tddim,font=\ttfamily\scriptsize]
        at ([xshift=\dimexpr\td@stripw+2mm\relax,yshift=13mm]current page.south west){#1};
    \else
      \node[rotate=90,anchor=west,inner xsep=1.2mm,inner ysep=0.6mm,
            fill=tdcard,text=tddim,font=\ttfamily\scriptsize]
        at ([xshift=\dimexpr-\td@stripw-2mm\relax,yshift=13mm]current page.south east){#1};
    \fi
  \end{tikzpicture}%
}
\makeatother

% A References frame that runs long is set [allowframebreaks]; beamer then
% appends a roman continuation counter to the title ("References II"), which
% reads like a second section. Blank it: both pages are just References.
\setbeamertemplate{frametitle continuation}{}

% ------------------------------------------------------------- statement
% The punchline slide, overridden deck level. The theme's own statement is
% letterspaced uppercase with a red uppercase second line; both lines shout and
% the red one competes with the first instead of supporting it.
%
% This is the Disciplined Vibe closing card's idea, editorial rather than
% centred: a short red rule, the claim flush left in large bold MIXED case with
% the red carried inline by \alertbf, then the consequence one step down in
% grey. The rule is the same motif as the subchapter divider, so dividers and
% punchlines read as one family, and flush left means a wrapping claim breaks
% cleanly instead of stranding an orphan word in the middle of the slide.
%
%   \begin{statement}
%     The algorithm is the \alertbf{small box in the middle}
%     \substatement{Everything else is this course.}
%   \end{statement}
%
% The theme's versions are deliberately shadowed, not patched: other decks share
% that file. \RenewDocumentEnvironment/Command so this fails loudly if the theme
% ever stops providing them.
\RenewDocumentEnvironment{statement}{+b}{%
  \par\vfill
  {\color{DHBWred}\rule{0.10\paperwidth}{2.5pt}}\par\vskip1.1em
  {\fontsize{28}{33}\selectfont\bfseries\raggedright #1\par}%
  \vfill\par}{}
% \substatement stays \protected so the surrounding grab cannot expand it.
\RenewDocumentCommand{\substatement}{m}{%
  \par\vskip0.9em{\large\color{tddim}\raggedright #1\par}}

% -------------------------------------------------------------- takeaway
% Also overridden deck level, and for the same reason as statement above: the
% theme's takeaway is \Large italic red centred, which shouts in a typeface
% that was never drawn to shout, and centring strands an orphan word when the
% line wraps. Here it joins the statement family instead: red rule, claim flush
% left in bold MIXED case, the red carried inline by \alertbf rather than
% painted over the whole sentence. One step smaller than statement, because a
% takeaway closes a chapter where a statement opens a question.
%
%   \begin{takeaway}
%     The gap between the two is the \alertbf{machine learning system}.
%   \end{takeaway}
\RenewDocumentEnvironment{takeaway}{+b}{%
  \par\vfill
  {\color{DHBWred}\rule{0.07\paperwidth}{2.5pt}}\par\vskip0.9em
  {\fontsize{22}{27}\selectfont\bfseries\raggedright #1\par}%
  \vfill\par}{}

% ---------------------------------------------------------- frame subtitle
% The theme's frametitle template prints the title and the hairline rule and
% drops \insertframesubtitle on the floor, so \framesubtitle{...} renders
% nothing. It is reinstated here, deck level, with the theme's own template
% copied verbatim and one branch added: grey, small, above the rule.
% Used for the source's own one line under a heading, e.g. the condition that
% names each of the five keyphrases in chapter 1.
\setbeamerfont{framesubtitle}{size=\small}
\setbeamercolor{framesubtitle}{fg=tddim}
\makeatletter
\setbeamertemplate{frametitle}{%
  \nointerlineskip%
  \vskip2.4ex%
  {\usebeamerfont{frametitle}\usebeamercolor[fg]{frametitle}\strut\tikz[remember picture,overlay]\coordinate (tdframetitle);\insertframetitle\par}%
  \ifx\insertframesubtitle\@empty\else
    \vskip0.35ex%
    {\usebeamerfont{framesubtitle}\usebeamercolor[fg]{framesubtitle}%
      \insertframesubtitle\par}%
  \fi
  \vskip0.2ex%
  {\color{tdrule}\rule{\textwidth}{0.4pt}}%
  \vskip0.4ex%
}
\makeatother

% Title on top of a side strip: \titleontop, called after \sideimagecover.
% The frametitle is typeset before the frame body, so a strip laid down in
% the body covers a long title's right end. The template above marks where
% the title starts with a zero-size coordinate; this draws the title again
% at that mark, over the photo.
\newcommand{\titleontop}{%
  \begin{tikzpicture}[remember picture,overlay]
    \node[anchor=base west,inner sep=0pt] at (tdframetitle)
      {\usebeamerfont{frametitle}\usebeamercolor[fg]{frametitle}\insertframetitle};
  \end{tikzpicture}%
}

% ------------------------------------------------------------- footer nav
% The theme's footer is beamer's mini-frame navigation: one hollow dot per
% frame, grouped under its section name. It is sized for a deck the shape of
% Disciplined Vibe, eight sections over fifty frames. This deck mirrors a
% seventeen block course: sixteen sections over about a hundred and eighty
% frames, which overruns the footer band by some 340pt on every single page.
% Neither knob fixes it, the dots alone are already three times the paper
% width. So the nav is replaced here, at deck level, with the same information
% compressed: the section you are in on the left, a progress track in the
% middle, the counter on the right. The theme is untouched and the other decks
% keep their dots. The track is drawn in tdfg on a tdrule ground rather than in
% red, so the footer never spends the slide's one red element.
% Drawn with \rule rather than tikz on purpose: a tikz coordinate resolves a
% bare length register through pgfmath, which hands back a unitless number that
% tikz then reads as centimetres, and the bar blows past the maximum dimension.
\newlength{\ieprogwd}
\newlength{\ieprogfill}
\setlength{\ieprogwd}{0.30\paperwidth}
% The fraction goes through pgfmath, which works in floating point. Doing it in
% TeX dimen arithmetic does not survive the first pass: before the aux file is
% written \inserttotalframenumber is stale, the \divide leaves \ieprogwd whole,
% and \ieprogwd times a frame number past about 120 exceeds the maximum TeX
% dimension. The clamp covers the same stale-aux case, where the current frame
% number can exceed the recorded total and the fill would run past its track.
\newcommand{\ieprogressbar}{%
  \ifnum\inserttotalframenumber>0\relax
    \ifnum\insertframenumber<\inserttotalframenumber\relax
      % The ratio is parenthesised so it is evaluated first. pgfmath keeps its
      % intermediates in a dimen register, so multiplying the track width by the
      % frame number before dividing overflows once past frame 120.
      \pgfmathsetlength{\ieprogfill}{\ieprogwd*(\insertframenumber/\inserttotalframenumber)}%
    \else
      \setlength{\ieprogfill}{\ieprogwd}%
    \fi
  \else
    \setlength{\ieprogfill}{0pt}%
  \fi
  \raisebox{0.6ex}{\makebox[\ieprogwd][l]{%
    \rlap{\textcolor{tdrule}{\rule{\ieprogwd}{1.1pt}}}%
    \textcolor{tdfg}{\rule{\ieprogfill}{1.1pt}}}}%
}
% The speaker's name leads the footer on every slide, so a photo taken during
% the talk carries it (Mark, 2026-10-26). The right edge watermark cannot do
% this: side strips and full page sketches cover it.
\newcommand{\iefootname}{Prof.\ Dr.-Ing.\ Mark Schutera\enspace\textperiodcentered\enspace}
\setbeamertemplate{footline}{%
  \begin{beamercolorbox}[wd=\paperwidth,ht=6mm,dp=5mm,leftskip=6mm,rightskip=6mm]{section in head/foot}%
    {\usebeamerfont{section in head/foot}\color{tddim}\iefootname\insertsectionhead}%
    \hfill
    \ieprogressbar
    \hfill
    {\usebeamerfont{section in head/foot}\color{tddim}%
      \insertframenumber\,/\,\inserttotalframenumber}%
  \end{beamercolorbox}%
}

%% (sharedAssets/beamer is on TEXINPUTS through each deck's .latexmkrc.)
%%
%% What it changes against the bare theme:
%%   footer       speaker name, section, progress track, counter; no dots
%%   \sidecaption up the strip's inner edge in typewriter grey, off the image
%%   statement, takeaway   red rule, flush left, bold mixed case, \alertbf
%%   \framesubtitle        printed, grey, above the title rule
%%   \titleontop           redraws a long title over a side strip
%%
%% COPIED VERBATIM from notes_intelligenceengineering/slides/
%% intelligence_engineering.tex (2026-11-20), which still carries its own
%% copy. The comments speak of "this deck" for that reason. Change both, or
%% move that deck onto this file.

\makeatletter
% \sidecaption shadowed for this deck. The theme sets the caption horizontal,
% white and italic INSIDE the strip, which a photo swallows; here it reads up
% the strip's outer edge in the placeholder's own typewriter grey, mirroring
% \slidesources on the opposite page edge, so a real image and its
% \phimgside placeholder caption the same way. 13mm clears the footer band.
% The node is filled with the page's own card colour and given a little
% horizontal padding, so wherever a long caption runs up past the frametitle's
% hairline rule it masks that rule instead of crossing it. The caption always
% sits OUTSIDE the image strip, on the card, so the fill is invisible; it only
% ever paints over the rule. Keep captions under ~50 characters anyway: the
% node has no text width and the page clips whatever runs past the top edge.
\renewcommand{\sidecaption}[1]{%
  \setlength{\td@stripw}{\td@sidefrac\paperwidth}%
  \def\td@westtest{west}%
  \begin{tikzpicture}[remember picture,overlay]
    \ifx\td@sideside\td@westtest
      \node[rotate=90,anchor=west,inner xsep=1.2mm,inner ysep=0.6mm,
            fill=tdcard,text=tddim,font=\ttfamily\scriptsize]
        at ([xshift=\dimexpr\td@stripw+2mm\relax,yshift=13mm]current page.south west){#1};
    \else
      \node[rotate=90,anchor=west,inner xsep=1.2mm,inner ysep=0.6mm,
            fill=tdcard,text=tddim,font=\ttfamily\scriptsize]
        at ([xshift=\dimexpr-\td@stripw-2mm\relax,yshift=13mm]current page.south east){#1};
    \fi
  \end{tikzpicture}%
}
\makeatother

% A References frame that runs long is set [allowframebreaks]; beamer then
% appends a roman continuation counter to the title ("References II"), which
% reads like a second section. Blank it: both pages are just References.
\setbeamertemplate{frametitle continuation}{}

% ------------------------------------------------------------- statement
% The punchline slide, overridden deck level. The theme's own statement is
% letterspaced uppercase with a red uppercase second line; both lines shout and
% the red one competes with the first instead of supporting it.
%
% This is the Disciplined Vibe closing card's idea, editorial rather than
% centred: a short red rule, the claim flush left in large bold MIXED case with
% the red carried inline by \alertbf, then the consequence one step down in
% grey. The rule is the same motif as the subchapter divider, so dividers and
% punchlines read as one family, and flush left means a wrapping claim breaks
% cleanly instead of stranding an orphan word in the middle of the slide.
%
%   \begin{statement}
%     The algorithm is the \alertbf{small box in the middle}
%     \substatement{Everything else is this course.}
%   \end{statement}
%
% The theme's versions are deliberately shadowed, not patched: other decks share
% that file. \RenewDocumentEnvironment/Command so this fails loudly if the theme
% ever stops providing them.
\RenewDocumentEnvironment{statement}{+b}{%
  \par\vfill
  {\color{DHBWred}\rule{0.10\paperwidth}{2.5pt}}\par\vskip1.1em
  {\fontsize{28}{33}\selectfont\bfseries\raggedright #1\par}%
  \vfill\par}{}
% \substatement stays \protected so the surrounding grab cannot expand it.
\RenewDocumentCommand{\substatement}{m}{%
  \par\vskip0.9em{\large\color{tddim}\raggedright #1\par}}

% -------------------------------------------------------------- takeaway
% Also overridden deck level, and for the same reason as statement above: the
% theme's takeaway is \Large italic red centred, which shouts in a typeface
% that was never drawn to shout, and centring strands an orphan word when the
% line wraps. Here it joins the statement family instead: red rule, claim flush
% left in bold MIXED case, the red carried inline by \alertbf rather than
% painted over the whole sentence. One step smaller than statement, because a
% takeaway closes a chapter where a statement opens a question.
%
%   \begin{takeaway}
%     The gap between the two is the \alertbf{machine learning system}.
%   \end{takeaway}
\RenewDocumentEnvironment{takeaway}{+b}{%
  \par\vfill
  {\color{DHBWred}\rule{0.07\paperwidth}{2.5pt}}\par\vskip0.9em
  {\fontsize{22}{27}\selectfont\bfseries\raggedright #1\par}%
  \vfill\par}{}

% ---------------------------------------------------------- frame subtitle
% The theme's frametitle template prints the title and the hairline rule and
% drops \insertframesubtitle on the floor, so \framesubtitle{...} renders
% nothing. It is reinstated here, deck level, with the theme's own template
% copied verbatim and one branch added: grey, small, above the rule.
% Used for the source's own one line under a heading, e.g. the condition that
% names each of the five keyphrases in chapter 1.
\setbeamerfont{framesubtitle}{size=\small}
\setbeamercolor{framesubtitle}{fg=tddim}
\makeatletter
\setbeamertemplate{frametitle}{%
  \nointerlineskip%
  \vskip2.4ex%
  {\usebeamerfont{frametitle}\usebeamercolor[fg]{frametitle}\strut\tikz[remember picture,overlay]\coordinate (tdframetitle);\insertframetitle\par}%
  \ifx\insertframesubtitle\@empty\else
    \vskip0.35ex%
    {\usebeamerfont{framesubtitle}\usebeamercolor[fg]{framesubtitle}%
      \insertframesubtitle\par}%
  \fi
  \vskip0.2ex%
  {\color{tdrule}\rule{\textwidth}{0.4pt}}%
  \vskip0.4ex%
}
\makeatother

% Title on top of a side strip: \titleontop, called after \sideimagecover.
% The frametitle is typeset before the frame body, so a strip laid down in
% the body covers a long title's right end. The template above marks where
% the title starts with a zero-size coordinate; this draws the title again
% at that mark, over the photo.
\newcommand{\titleontop}{%
  \begin{tikzpicture}[remember picture,overlay]
    \node[anchor=base west,inner sep=0pt] at (tdframetitle)
      {\usebeamerfont{frametitle}\usebeamercolor[fg]{frametitle}\insertframetitle};
  \end{tikzpicture}%
}

% ------------------------------------------------------------- footer nav
% The theme's footer is beamer's mini-frame navigation: one hollow dot per
% frame, grouped under its section name. It is sized for a deck the shape of
% Disciplined Vibe, eight sections over fifty frames. This deck mirrors a
% seventeen block course: sixteen sections over about a hundred and eighty
% frames, which overruns the footer band by some 340pt on every single page.
% Neither knob fixes it, the dots alone are already three times the paper
% width. So the nav is replaced here, at deck level, with the same information
% compressed: the section you are in on the left, a progress track in the
% middle, the counter on the right. The theme is untouched and the other decks
% keep their dots. The track is drawn in tdfg on a tdrule ground rather than in
% red, so the footer never spends the slide's one red element.
% Drawn with \rule rather than tikz on purpose: a tikz coordinate resolves a
% bare length register through pgfmath, which hands back a unitless number that
% tikz then reads as centimetres, and the bar blows past the maximum dimension.
\newlength{\ieprogwd}
\newlength{\ieprogfill}
\setlength{\ieprogwd}{0.30\paperwidth}
% The fraction goes through pgfmath, which works in floating point. Doing it in
% TeX dimen arithmetic does not survive the first pass: before the aux file is
% written \inserttotalframenumber is stale, the \divide leaves \ieprogwd whole,
% and \ieprogwd times a frame number past about 120 exceeds the maximum TeX
% dimension. The clamp covers the same stale-aux case, where the current frame
% number can exceed the recorded total and the fill would run past its track.
\newcommand{\ieprogressbar}{%
  \ifnum\inserttotalframenumber>0\relax
    \ifnum\insertframenumber<\inserttotalframenumber\relax
      % The ratio is parenthesised so it is evaluated first. pgfmath keeps its
      % intermediates in a dimen register, so multiplying the track width by the
      % frame number before dividing overflows once past frame 120.
      \pgfmathsetlength{\ieprogfill}{\ieprogwd*(\insertframenumber/\inserttotalframenumber)}%
    \else
      \setlength{\ieprogfill}{\ieprogwd}%
    \fi
  \else
    \setlength{\ieprogfill}{0pt}%
  \fi
  \raisebox{0.6ex}{\makebox[\ieprogwd][l]{%
    \rlap{\textcolor{tdrule}{\rule{\ieprogwd}{1.1pt}}}%
    \textcolor{tdfg}{\rule{\ieprogfill}{1.1pt}}}}%
}
% The speaker's name leads the footer on every slide, so a photo taken during
% the talk carries it (Mark, 2026-10-26). The right edge watermark cannot do
% this: side strips and full page sketches cover it.
\newcommand{\iefootname}{Prof.\ Dr.-Ing.\ Mark Schutera\enspace\textperiodcentered\enspace}
\setbeamertemplate{footline}{%
  \begin{beamercolorbox}[wd=\paperwidth,ht=6mm,dp=5mm,leftskip=6mm,rightskip=6mm]{section in head/foot}%
    {\usebeamerfont{section in head/foot}\color{tddim}\iefootname\insertsectionhead}%
    \hfill
    \ieprogressbar
    \hfill
    {\usebeamerfont{section in head/foot}\color{tddim}%
      \insertframenumber\,/\,\inserttotalframenumber}%
  \end{beamercolorbox}%
}
